\documentclass[11pt]{article}
\usepackage{graphicx,amsmath,amssymb,amsthm,booktabs,multirow,xcolor,hyperref,geometry,longtable,array}
\usepackage[numbers,sort&compress,super]{natbib}
\geometry{margin=1in}

\title{Supplementary Information:\\Physics World Models for Computational Imaging}
\author{Chengshuai Yang and Xin Yuan}
\date{}

\newtheorem{theorem}{Theorem}
\newtheorem{proposition}[theorem]{Proposition}

\begin{document}
\maketitle

\renewcommand{\thetable}{S\arabic{table}}
\renewcommand{\thefigure}{S\arabic{figure}}
\renewcommand{\theequation}{S\arabic{equation}}
\setcounter{table}{0}
\setcounter{figure}{0}
\setcounter{equation}{0}

\tableofcontents
\newpage

%% ============================================================
%% SUPPLEMENTARY NOTE 1
%% ============================================================
\section{Supplementary Note 1: Triad Decomposition Mathematical Derivations}
\label{sec:triad}

The Triad Decomposition formalizes three successive gates that every computational-imaging
reconstruction must pass.  We derive rigorous bounds for each gate and show how
the \emph{recovery ratio} $\rho$ summarizes the overall reconstruction quality.

\subsection{Gate 1 --- Information-Theoretic Limit (Compression)}

Let the forward operator $\mathbf{H}\in\mathbb{R}^{m\times n}$ map an
$n$-dimensional scene $\mathbf{x}$ to $m$ measurements $\mathbf{y}$.
Define the \emph{compression ratio} $\gamma = m/n$.

\begin{theorem}[Compression bound]
For any estimator $\hat{\mathbf{x}}(\mathbf{y})$, the minimum achievable
mean-squared error satisfies
\begin{equation}
  \mathrm{MSE}_{\min} \;\geq\;
  \frac{1}{n}\sum_{i=1}^{n-m} \sigma_i^2(\mathbf{x}),
  \label{eq:gate1}
\end{equation}
where $\sigma_i^2(\mathbf{x})$ denotes the variance of $\mathbf{x}$ along the
$i$-th null-space direction of $\mathbf{H}$.
\end{theorem}

\noindent\textbf{Interpretation.}
The null space of $\mathbf{H}$ has dimension $\dim\mathcal{N}(\mathbf{H}) =
n - \mathrm{rank}(\mathbf{H}) \geq n - m$.  Information lost in the null space
cannot be recovered without a prior; Gate~1 therefore sets an
\emph{information-theoretic floor} on PSNR:
\begin{equation}
  \mathrm{PSNR}_{\max}^{(G1)}
  = 10\log_{10}\!\!\left(\frac{\|\mathbf{x}\|_\infty^2}{\mathrm{MSE}_{\min}}\right).
\end{equation}
Modalities with higher compression ratios (e.g.\ CACTI with $B$
frames compressed into one) face a proportionally lower ceiling.

\subsection{Gate 2 --- SNR-Dependent PSNR Bound (Noise)}

Given additive noise $\mathbf{n}\sim\mathcal{N}(0,\sigma_n^2\mathbf{I})$,
the measurement SNR is $\mathrm{SNR} = \|\mathbf{H}\mathbf{x}\|^2 /
(m\,\sigma_n^2)$.

\begin{theorem}[Noise bound]
Under matched forward model and Gaussian noise, the best achievable PSNR is
\begin{equation}
  \mathrm{PSNR}_{\max}^{(G2)}
  = 10\log_{10}(\mathrm{SNR}) + C_{\mathcal{M}},
  \label{eq:gate2}
\end{equation}
where $C_{\mathcal{M}}$ is a modality-dependent constant that absorbs
the conditioning of $\mathbf{H}$ and the prior strength:
\begin{equation}
  C_{\mathcal{M}}
  = 10\log_{10}\!\left(
      \frac{n\,\|\mathbf{x}\|_\infty^2}{\|\mathbf{x}\|^2}
      \cdot \kappa^{-2}(\mathbf{H})
    \right),
\end{equation}
with $\kappa(\mathbf{H})$ the condition number of $\mathbf{H}$ restricted to
its column space.
\end{theorem}

\noindent\textbf{Interpretation.}
Gate~2 is a \emph{noise ceiling}: no algorithm can exceed
$\mathrm{PSNR}_{\max}^{(G2)}$ regardless of computational budget.
Empirically, modalities such as MRI ($C_{\mathrm{MRI}}\approx 8$\,dB) are more
noise-tolerant than CASSI ($C_{\mathrm{CASSI}}\approx 3$\,dB) due to
favorable operator conditioning.

\subsection{Gate 3 --- Sensitivity to Model Mismatch (Calibration)}

Let $\boldsymbol{\theta}\in\mathbb{R}^p$ collect the calibration parameters
(e.g.\ mask shift, PSF blur, coil sensitivity) and let $\mathbf{H}_{\boldsymbol{\theta}}$
denote the parameterized forward operator.

\begin{theorem}[Calibration sensitivity]
For small perturbation $\delta\boldsymbol{\theta}$ around the true parameter
$\boldsymbol{\theta}^*$, the PSNR degradation is
\begin{equation}
  \Delta\mathrm{PSNR}
  \approx -\frac{10}{\ln 10}\,
  \frac{\delta\boldsymbol{\theta}^\top
        \mathbf{J}^\top \mathbf{J}\,
        \delta\boldsymbol{\theta}}
       {\mathrm{MSE}_0},
  \label{eq:gate3}
\end{equation}
where $\mathbf{J} = \partial(\mathbf{H}_{\boldsymbol{\theta}}\mathbf{x})/
\partial\boldsymbol{\theta}\big|_{\boldsymbol{\theta}^*}$ is the parameter
Jacobian and $\mathrm{MSE}_0$ is the noise-only MSE at
$\boldsymbol{\theta}^*$.
\end{theorem}

\noindent\textbf{Per-parameter sensitivity.}
Restricting to a single parameter $\theta_j$:
\begin{equation}
  \frac{d\,\mathrm{PSNR}}{d\,\theta_j}
  = -\frac{10}{\ln 10}\,
    \frac{\|\mathbf{J}_{:,j}\|^2}{\mathrm{MSE}_0}\,\delta\theta_j
  + \mathcal{O}(\delta\theta_j^2).
\end{equation}
The first-order Taylor expansion shows that PSNR degrades linearly in
$|\delta\theta_j|$ for small mismatches and quadratically for larger ones.

\subsection{Recovery Ratio}

We define the \emph{recovery ratio} as
\begin{equation}
  \rho
  = \frac{\mathrm{PSNR}_{\mathrm{achieved}} - \mathrm{PSNR}_{\mathrm{mismatch}}}
         {\mathrm{PSNR}_{\mathrm{ideal}} - \mathrm{PSNR}_{\mathrm{mismatch}}},
  \qquad 0 \leq \rho \leq 1.
  \label{eq:rho}
\end{equation}

\noindent In rare cases where Scenario~III exceeds Scenario~I (e.g., due to
regularization benefits from the corrected operator), $\rho$ may exceed~1.

\begin{proposition}
Under convex reconstruction loss with matched regularization and convex constraint on
$\boldsymbol{\theta}$, the recovery ratio satisfies $0 \leq \rho \leq 1$
for any estimator.
\end{proposition}

\begin{proof}[Proof sketch]
The oracle PSNR is the global optimum of the jointly convex problem in
$(\mathbf{x}, \boldsymbol{\theta})$.  The mismatch PSNR uses a fixed
$\boldsymbol{\theta}_0\neq\boldsymbol{\theta}^*$, yielding a suboptimal
feasible point.  Any correction step moves towards the optimum, so
$\mathrm{PSNR}_{\mathrm{mismatch}} \leq \mathrm{PSNR}_{\mathrm{achieved}}
\leq \mathrm{PSNR}_{\mathrm{oracle}}$, giving $\rho\in[0,1]$.
\end{proof}

\noindent Under the conditions of Proposition~1, $0 \leq \rho \leq 1$. In practice, beneficial regularization bias from the corrected operator can yield $\rho > 1$, as observed for CACTI.

\newpage

%% ============================================================
%% SUPPLEMENTARY NOTE 2
%% ============================================================
\section{Supplementary Note 2: Complete OperatorGraph Specification}
\label{sec:opgraph}

The \texttt{OperatorGraph} is the central abstraction of the PWM framework.
It represents a computational imaging pipeline as a directed acyclic graph (DAG)
of differentiable operators.

\subsection{Node Interface}

Every node $v$ in the graph implements the following interface:

\begin{center}
\renewcommand{\arraystretch}{1.3}
\begin{tabular}{lp{8cm}}
\toprule
\textbf{Method} & \textbf{Contract} \\
\midrule
\texttt{forward(x) $\to$ y} &
  Apply the physical forward model.  Input shape: \texttt{(B, C, *spatial)}.
  Output shape determined by the operator (e.g.\ compression changes spatial dims). \\
\texttt{adjoint(y) $\to$ x} &
  Apply the adjoint (transpose) operator.  Must satisfy
  $\langle \mathbf{H}\mathbf{x}, \mathbf{y}\rangle =
   \langle \mathbf{x}, \mathbf{H}^\top\mathbf{y}\rangle$
  to numerical precision ($<10^{-6}$ relative error). \\
\texttt{shape\_in} / \texttt{shape\_out} &
  Static shape annotations for compile-time validation. \\
\texttt{dtype} &
  Data type contract: \texttt{float32} (default) or \texttt{complex64}. \\
\texttt{parameters() $\to$ dict} &
  Returns learnable calibration parameters with their current values. \\
\bottomrule
\end{tabular}
\end{center}

\subsection{Edge Semantics}

An edge $e = (u, v)$ represents data flow: the output tensor of node $u$ is
fed as input to node $v$.  Edges carry the following metadata:
\begin{itemize}
  \item \textbf{Tensor shape}: inferred at compile time via shape propagation.
  \item \textbf{Data type}: must match between source output and target input.
  \item \textbf{Optional transform}: lightweight reshape or permute operations
        (e.g.\ channel stacking for multi-frame CACTI data).
\end{itemize}

\subsection{Compilation Algorithm}

Given an \texttt{OperatorGraph} $G = (V, E)$:

\begin{enumerate}
  \item \textbf{Topological sort.}  Compute a valid execution order
        $v_1, v_2, \ldots, v_{|V|}$ using Kahn's algorithm.
        Raise \texttt{CycleError} if the graph contains a cycle.
  \item \textbf{Shape propagation.}  Starting from the source node(s), propagate
        tensor shapes through each operator's \texttt{shape\_out} to validate
        all edge contracts.
  \item \textbf{Automatic adjoint chain.}  Construct the adjoint graph
        $G^\top$ by reversing all edges and replacing each node's
        \texttt{forward} with \texttt{adjoint}.  The adjoint graph is used for
        gradient-based reconstruction and the Triad Decomposition sensitivity analysis.
  \item \textbf{Fusion (optional).}  Consecutive linear operators are fused
        into a single matrix--vector product when dimensions permit, reducing
        memory traffic.
\end{enumerate}

\subsection{JSON Serialization Schema}

The graph is serialized to a JSON document for reproducibility and sharing:

\begin{verbatim}
{
  "version": "0.3.0",
  "nodes": [
    {
      "id": "mask",
      "type": "CassiMask",
      "params": {"shift_r": 0.0, "shift_c": 0.0},
      "shape_in": [1, 28, 256, 256],
      "shape_out": [1, 1, 256, 310]
    }
  ],
  "edges": [
    {"src": "mask", "dst": "detector", "dtype": "float32"}
  ],
  "metadata": {
    "modality": "CASSI",
    "created": "2026-02-01T00:00:00Z"
  }
}
\end{verbatim}

\newpage

%% ============================================================
%% SUPPLEMENTARY NOTE 3
%% ============================================================
\section{Supplementary Table S1: All 16 Correction Configurations}
\label{sec:params}

Table~\ref{tab:correction_params} reports the correction results for all 16
correction configurations across the validated modalities.  The first 9 modalities
have full end-to-end validation; the remaining 7 are designated Phase~2/4
additions with partial or planned validation.

\begin{longtable}{>{\raggedright}p{2.0cm} >{\raggedright}p{2.8cm}
                  r r r r}
\caption{Complete correction parameter summary.  $\Delta$PSNR is the
improvement from Scenario~II (mismatch) to Scenario~III (corrected).
CASSI correction results in this table use GAP-TV as the reconstruction solver; the main-text deep dive reports results for GAP-TV, MST-L, and HDNet.
Scenario~I (ideal) PSNR values and recovery ratios ($\rho$) for all configurations are available in the RunBundle manifests in the code repository.}
\label{tab:correction_params} \\
\toprule
\textbf{Modality} & \textbf{Mismatch Parameter} &
  \textbf{Sc.~II} & \textbf{Sc.~III} &
  \textbf{$\Delta$PSNR} & \textbf{RMSE} \\
  & & \textbf{(dB)} & \textbf{(dB)} & \textbf{(dB)} & \\
\midrule
\endfirsthead
\toprule
\textbf{Modality} & \textbf{Mismatch Parameter} &
  \textbf{Sc.~II} & \textbf{Sc.~III} &
  \textbf{$\Delta$PSNR} & \textbf{RMSE} \\
  & & \textbf{(dB)} & \textbf{(dB)} & \textbf{(dB)} & \\
\midrule
\endhead
\midrule
\multicolumn{6}{r}{\emph{Continued on next page}} \\
\endfoot
\bottomrule
\endlastfoot
%% --- Validated modalities (9) ---
Matrix\footnote{Matrix (Generic) is a canonical test configuration that uses the same gain-bias operator template as SPC; identical results confirm pipeline consistency rather than representing an independent modality.}
  & gain\_bias
  & 11.14 & 23.35 & +12.21
  & 0.0042 \\
CT
  & center\_of\_rotation
  & 13.41 & 24.09 & +10.68
  & 0.0031 \\
CACTI
  & mask\_timing
  & 14.48 & 37.42 & +22.94
  & 0.0018 \\
Lensless
  & psf\_shift
  & 23.48 & 27.03 & +3.55
  & 0.0089 \\
MRI
  & coil\_sensitivities
  & 6.94 & 55.19 & +48.25
  & 0.0003 \\
SPC
  & gain\_bias
  & 11.14 & 23.35 & +12.21
  & 0.0042 \\
CASSI (Alg~1)\footnote{Reported values are averaged over 10 test scenes using GAP-TV as the reconstruction solver. Sc.~II from InverseNet validation; Sc.~III values are oracle upper bounds (true operator applied to mismatched data).}
  & mask\_geo + dispersion
  & 20.96 & 21.50 & +0.54
  & 0.0105 \\
CASSI (Alg~2)\footnote{Note: the correction-table PSNR reflects GAP-TV reconstruction; state-of-the-art solvers (e.g., MST-L) achieve 27.33\,dB oracle PSNR under the corrected operator (see main text).}
  & mask\_geo + dispersion
  & 20.96 & 21.72 & +0.76
  & 0.0098 \\
Ptychography
  & position\_offset
  & 17.35 & 24.44 & +7.09
  & 0.0063 \\
\midrule
%% --- Phase 2/4 additions (7) ---
\multicolumn{6}{l}{\emph{Phase 2/4 additions --- validation in progress}} \\
\midrule
Holography
  & propagation\_distance
  & --- & --- & ---
  & --- \\
FPM
  & LED\_positions
  & --- & --- & ---
  & --- \\
Phase Retrieval
  & support\_constraint
  & --- & --- & ---
  & --- \\
OCT
  & dispersion\_coeffs
  & --- & --- & ---
  & --- \\
Fluorescence
  & PSF\_aberration
  & --- & --- & ---
  & --- \\
Ultrasound
  & speed\_of\_sound
  & --- & --- & ---
  & --- \\
Radar/SAR
  & motion\_compensation
  & --- & --- & ---
  & --- \\
\end{longtable}

\newpage

%% ============================================================
%% SUPPLEMENTARY TABLE 1
%% ============================================================
\section{Supplementary Table S2: CASSI Per-Scene Results}
\label{sec:cassi_scenes}

Table~\ref{tab:cassi_scenes} reports PSNR (dB) for 10 KAIST scenes across
five representative methods under four scenarios.\footnote{This supplementary table reports five methods (GAP-TV, TwIST, DeSCI, MST, HDNet) to provide comprehensive coverage; the main text focuses on three methods (GAP-TV, MST-L, HDNet) for clarity. The supplementary table uses MST (base variant); the main text reports MST-L (large variant), which has higher Scenario~I performance but similar Scenario~II collapse.}

\begin{footnotesize}
\begin{longtable}{l l r r r r}
\caption{CASSI per-scene PSNR (dB).  Sc.~I = ideal forward model;
Sc.~II = mismatched (baseline); Sc.~III-A1 = correction Algorithm~1; Verify = pipeline sanity check (ideal-condition data with ideal operator, confirming pipeline integrity to $\pm 0.01$\,dB from Sc.~I).
GAP-TV values for Sc.~I, Sc.~II, and Sc.~III-A1 are from InverseNet validation~\cite{yang2026inversenet} under the 5-parameter mismatch ($dx{=}0.5$\,px, $dy{=}0.3$\,px, $\theta{=}0.1^\circ$, $a_1{=}2.02$, $\alpha{=}0.15^\circ$). Sc.~III-A1 entries marked $^*$ use Oracle Mask values (true operator applied to mismatched data) as correction upper bounds. Algorithm~2 results are reported in Table~\ref{tab:correction_params}.}
\label{tab:cassi_scenes} \\
\toprule
\textbf{Scene} & \textbf{Method} &
  \textbf{Sc.~I} & \textbf{Sc.~II} & \textbf{Sc.~III-A1} & \textbf{Verify} \\
\midrule
\endfirsthead
\toprule
\textbf{Scene} & \textbf{Method} &
  \textbf{Sc.~I} & \textbf{Sc.~II} & \textbf{Sc.~III-A1} & \textbf{Verify} \\
\midrule
\endhead
\bottomrule
\endlastfoot
%% Scene 1
scene01 & GAP-TV        & 26.49 & 24.08 & 24.18$^*$ & 26.49 \\
scene01 & TwIST         & 27.18 & 15.79 & 21.34 & 27.18 \\
scene01 & DeSCI         & 28.01 & 14.92 & 20.15 & 28.01 \\
scene01 & MST           & 33.47 & 16.84 & 22.61 & 33.47 \\
scene01 & HDNet         & 34.12 & 17.02 & 23.10 & 34.12 \\
\midrule
%% Scene 2
scene02 & GAP-TV        & 24.60 & 21.89 & 22.82$^*$ & 24.60 \\
scene02 & TwIST         & 25.73 & 15.21 & 20.48 & 25.73 \\
scene02 & DeSCI         & 26.95 & 14.38 & 19.54 & 26.95 \\
scene02 & MST           & 31.84 & 16.12 & 21.73 & 31.84 \\
scene02 & HDNet         & 32.56 & 16.45 & 22.01 & 32.56 \\
\midrule
%% Scene 3
scene03 & GAP-TV        & 25.96 & 18.62 & 19.68$^*$ & 25.96 \\
scene03 & TwIST         & 28.04 & 16.35 & 21.87 & 28.04 \\
scene03 & DeSCI         & 29.22 & 15.64 & 20.93 & 29.22 \\
scene03 & MST           & 34.67 & 17.48 & 23.25 & 34.67 \\
scene03 & HDNet         & 35.31 & 17.71 & 23.68 & 35.31 \\
\midrule
%% Scene 4
scene04 & GAP-TV        & 28.36 & 23.37 & 24.41$^*$ & 28.36 \\
scene04 & TwIST         & 31.09 & 18.02 & 23.75 & 31.09 \\
scene04 & DeSCI         & 32.14 & 17.11 & 22.64 & 32.14 \\
scene04 & MST           & 37.82 & 19.33 & 25.47 & 37.82 \\
scene04 & HDNet         & 38.45 & 19.68 & 25.94 & 38.45 \\
\midrule
%% Scene 5
scene05 & GAP-TV        & 23.66 & 20.39 & 21.27$^*$ & 23.66 \\
scene05 & TwIST         & 26.54 & 15.53 & 20.84 & 26.54 \\
scene05 & DeSCI         & 27.68 & 14.72 & 19.87 & 27.68 \\
scene05 & MST           & 32.91 & 16.67 & 22.14 & 32.91 \\
scene05 & HDNet         & 33.58 & 16.94 & 22.53 & 33.58 \\
\midrule
%% Scene 6
scene06 & GAP-TV        & 22.34 & 20.39 & 21.00$^*$ & 22.34 \\
scene06 & TwIST         & 29.18 & 16.74 & 22.31 & 29.18 \\
scene06 & DeSCI         & 30.42 & 15.97 & 21.24 & 30.42 \\
scene06 & MST           & 35.89 & 18.05 & 23.87 & 35.89 \\
scene06 & HDNet         & 36.54 & 18.32 & 24.25 & 36.54 \\
\midrule
%% Scene 7
scene07 & GAP-TV        & 23.51 & 20.56 & 21.07$^*$ & 23.51 \\
scene07 & TwIST         & 24.76 & 14.32 & 19.64 & 24.76 \\
scene07 & DeSCI         & 25.84 & 13.55 & 18.73 & 25.84 \\
scene07 & MST           & 30.72 & 15.48 & 21.02 & 30.72 \\
scene07 & HDNet         & 31.35 & 15.73 & 21.38 & 31.35 \\
\midrule
%% Scene 8
scene08 & GAP-TV        & 22.16 & 20.57 & 21.10$^*$ & 22.16 \\
scene08 & TwIST         & 28.53 & 16.52 & 22.04 & 28.53 \\
scene08 & DeSCI         & 29.71 & 15.78 & 21.12 & 29.71 \\
scene08 & MST           & 35.14 & 17.83 & 23.55 & 35.14 \\
scene08 & HDNet         & 35.79 & 18.08 & 23.94 & 35.79 \\
\midrule
%% Scene 9
scene09 & GAP-TV        & 23.03 & 19.14 & 20.61$^*$ & 23.03 \\
scene09 & TwIST         & 26.89 & 15.64 & 21.02 & 26.89 \\
scene09 & DeSCI         & 28.04 & 14.91 & 20.08 & 28.04 \\
scene09 & MST           & 33.28 & 16.87 & 22.35 & 33.28 \\
scene09 & HDNet         & 33.94 & 17.12 & 22.71 & 33.94 \\
\midrule
%% Scene 10
scene10 & GAP-TV        & 23.27 & 20.58 & 21.04$^*$ & 23.27 \\
scene10 & TwIST         & 30.34 & 17.54 & 23.21 & 30.34 \\
scene10 & DeSCI         & 31.52 & 16.73 & 22.14 & 31.52 \\
scene10 & MST           & 37.05 & 18.89 & 24.91 & 37.05 \\
scene10 & HDNet         & 37.68 & 19.15 & 25.34 & 37.68 \\
\end{longtable}
\end{footnotesize}

\newpage

%% ============================================================
%% SUPPLEMENTARY TABLE 2
%% ============================================================
\section{Supplementary Table S3: Full 26-Modality Benchmark}
\label{sec:26mod}

Table~\ref{tab:26mod} lists all 26 modalities in the PWM benchmark, grouped
by physical carrier tier.

\begin{longtable}{r l l r r l l}
\caption{Full 26-modality benchmark.  PSNR values are representative
(Scenario~I, default solver).  ``Ref.\ PSNR'' is the best published result
where available.}
\label{tab:26mod} \\
\toprule
\textbf{\#} & \textbf{Modality} & \textbf{Domain} &
  \textbf{PSNR} & \textbf{Ref.\ PSNR} & \textbf{Status} & \textbf{Solver} \\
  & & & \textbf{(dB)} & \textbf{(dB)} & & \\
\midrule
\endfirsthead
\toprule
\textbf{\#} & \textbf{Modality} & \textbf{Domain} &
  \textbf{PSNR} & \textbf{Ref.\ PSNR} & \textbf{Status} & \textbf{Solver} \\
  & & & \textbf{(dB)} & \textbf{(dB)} & & \\
\midrule
\endhead
\bottomrule
\endlastfoot
%% Group 1: Incoherent light
1  & Matrix (toy)         & Incoherent     & 23.35 & ---   & Validated & ADMM-TV \\
2  & SPC                  & Incoherent     & 23.35 & 24.1  & Validated & ADMM-TV \\
3  & CASSI                & Incoherent     & 24.34 & 34.2  & Validated & GAP-TV \\
4  & CACTI                & Incoherent     & 35.33 & 33.4  & Validated & EfficientSCI \\
5  & Lensless             & Incoherent     & 27.03 & 28.5  & Validated & ADMM-TV \\
6  & Fluorescence         & Incoherent     & ---   & 30.1  & Phase 4   & --- \\
%% Group 2: Coherent light
7  & Holography           & Coherent       & ---   & 32.7  & Phase 2   & --- \\
8  & Ptychography         & Coherent       & 24.44 & 26.8  & Validated & ePIE \\
9  & FPM                  & Coherent       & ---   & 31.5  & Phase 2   & --- \\
10 & Phase Retrieval      & Coherent       & ---   & 29.3  & Phase 2   & --- \\
11 & CDI                  & Coherent       & ---   & 27.4  & Phase 4   & --- \\
12 & OCT                  & Coherent       & ---   & 35.2  & Phase 4   & --- \\
%% Group 3: X-ray / electron
13 & CT (parallel)        & X-ray          & 24.09 & 42.1  & Validated & FBP + TV \\
14 & CT (fan-beam)        & X-ray          & ---   & 40.8  & Phase 2   & --- \\
15 & CT (cone-beam)       & X-ray          & ---   & 39.5  & Phase 2   & --- \\
16 & Electron Pty.        & Electron       & ---   & 25.9  & Phase 4   & --- \\
%% Group 4: RF / acoustic
17 & MRI (Cartesian)      & RF             & 55.19 & 42.3  & Validated & CG-SENSE \\
18 & MRI (radial)         & RF             & ---   & 38.7  & Phase 2   & --- \\
19 & MRI (spiral)         & RF             & ---   & 37.2  & Phase 4   & --- \\
20 & Ultrasound           & Acoustic       & ---   & 31.8  & Phase 4   & --- \\
21 & Radar / SAR          & RF             & ---   & 28.4  & Phase 4   & --- \\
%% Group 5: Computational / hybrid
22 & NeRF                 & Multi-view     & ---   & 31.7  & Phase 2   & --- \\
23 & 3DGS                 & Multi-view     & ---   & 33.2  & Phase 2   & --- \\
24 & Light Field          & Multi-view     & ---   & 36.1  & Phase 4   & --- \\
25 & Structured Illum.    & Incoherent     & ---   & 32.4  & Phase 4   & --- \\
26 & Ghost Imaging        & Incoherent     & ---   & 22.8  & Phase 4   & --- \\
\end{longtable}

\newpage

%% ============================================================
%% SUPPLEMENTARY TABLE 3
%% ============================================================
\section{Supplementary Table S4: YAML Registry Summary}
\label{sec:yaml}

Table~\ref{tab:yaml} summarizes the nine YAML registry files (totalling
7,034 lines) that define the PWM framework's modular architecture.

\begin{table}[htbp]
\centering
\caption{YAML registry files in \texttt{packages/pwm\_core/contrib/}.}
\label{tab:yaml}
\begin{tabular}{r l r p{6.5cm}}
\toprule
\textbf{\#} & \textbf{File} & \textbf{Lines} & \textbf{Purpose} \\
\midrule
1 & \texttt{modalities.yaml}           & 1200 & Defines 64 modality entries with keywords, forward model equations, and default solvers. \\
2 & \texttt{graph\_templates.yaml}     & 1400 & OperatorGraph skeletons for 89 templates across 64 modalities. \\
3 & \texttt{photon\_db.yaml}           & 624 & Photon models parameterized by source power, quantum efficiency, exposure, and detector. \\
4 & \texttt{mismatch\_db.yaml}         & 797 & Mismatch parameters, perturbation ranges, and correction methods per modality. \\
5 & \texttt{compression\_db.yaml}      & 1186 & Recoverability tables mapping compression ratio to expected PSNR with provenance. \\
6 & \texttt{solver\_registry.yaml}     & 650 & Maps solver names to implementations (ADMM~\cite{boyd2011admm}, FISTA~\cite{beck2009fista}, GAP-TV, PnP, learned). \\
7 & \texttt{primitives.yaml}           & 450 & Primitive operator metadata (type, adjoint availability, differentiability). \\
8 & \texttt{dataset\_registry.yaml}    & 380 & Links modalities to benchmark datasets (KAIST, fastMRI, etc.). \\
9 & \texttt{acceptance\_thresholds.yaml} & 347 & Pass/fail thresholds per metric for automated validation. \\
\bottomrule
\end{tabular}
\end{table}

\newpage

%% ============================================================
%% SUPPLEMENTARY NOTE 4
%% ============================================================
\section{Supplementary Note 4: RunBundle Schema}
\label{sec:runbundle}

The \texttt{RunBundle} is PWM's reproducibility artifact.  Every experiment
produces a RunBundle that captures the complete computational state.

\subsection{RunBundle v0.3.0 Specification}

\begin{longtable}{l l p{7cm}}
\caption{RunBundle v0.3.0 field specification.}
\label{tab:runbundle} \\
\toprule
\textbf{Field} & \textbf{Type} & \textbf{Description} \\
\midrule
\endfirsthead
\toprule
\textbf{Field} & \textbf{Type} & \textbf{Description} \\
\midrule
\endhead
\bottomrule
\endlastfoot
\texttt{version}        & string   & Schema version, currently \texttt{"0.3.0"}. \\
\texttt{git\_hash}      & string   & Full 40-character SHA-1 of the commit used to generate results. \\
\texttt{rng\_seeds}     & object   & Random seeds for Python (\texttt{random}, \texttt{numpy}), PyTorch (CPU, CUDA), and any solver-specific RNG. \\
\texttt{platform\_info} & object   & OS, Python version, GPU model, CUDA version, PyTorch version, and \texttt{pip freeze} hash. \\
\texttt{sha256\_hashes} & object   & SHA-256 checksums of all input data files (measurements, masks, ground truth). \\
\texttt{metrics}        & object   & Computed quality metrics: PSNR, SSIM, LPIPS, and any modality-specific metrics. \\
\texttt{timestamps}     & object   & ISO-8601 timestamps for start, end, and per-stage completion. \\
\texttt{operator\_state}& object   & Serialized OperatorGraph including all node parameters and edge metadata. \\
\texttt{triad\_report}  & object   & Triad Decomposition evaluation: gate pass/fail status, PSNR bounds, and recovery ratio $\rho$. \\
\texttt{correction\_trajectory} & array &
  Time series of correction iterations: $[\{$\texttt{iter}, \texttt{params},
  \texttt{loss}, \texttt{psnr}, \texttt{grad\_norm}$\}, \ldots]$.
  Enables convergence analysis and early-stopping diagnostics. \\
\texttt{solver\_config} & object   & Complete solver hyperparameters (algorithm, iterations, step size, regularization weight, denoiser checkpoint). \\
\texttt{modality}       & string   & Modality identifier matching \texttt{modalities.yaml}. \\
\texttt{scenario}       & string   & One of \texttt{I} (ideal), \texttt{II} (mismatch), \texttt{III} (corrected), \texttt{IV} (oracle\_mask). \\
\texttt{notes}          & string   & Free-form text field for experiment annotations. \\
\end{longtable}

\subsection{Integrity Verification}

A RunBundle can be verified by:
\begin{enumerate}
  \item Checking \texttt{git\_hash} matches the repository state.
  \item Re-computing \texttt{sha256\_hashes} from the data files.
  \item Re-running the solver with the stored \texttt{rng\_seeds} and
        \texttt{solver\_config} and comparing metrics within tolerance
        ($\pm0.01$\,dB PSNR).
\end{enumerate}

\newpage

%% ============================================================
%% SUPPLEMENTARY NOTE 5
%% ============================================================
\section{Supplementary Note 5: Computational Cost Analysis}
\label{sec:cost}

\subsection{Runtime per Modality}

All timings were measured on a single NVIDIA RTX 4090 GPU (24\,GB VRAM)
with PyTorch 2.1 and CUDA 12.1.

\begin{table}[htbp]
\centering
\caption{Computational cost for correction (Scenario~II $\to$ III).
RoIC = dB recovered per GPU-hour.}
\label{tab:cost}
\begin{tabular}{l r r r r}
\toprule
\textbf{Modality} & \textbf{Iterations} & \textbf{Runtime (s)} &
  \textbf{$\Delta$PSNR (dB)} & \textbf{RoIC (dB/GPU-hr)} \\
\midrule
Matrix           &   200 &     12 & +12.21 & 3663.0 \\
CT               &   500 &     85 & +10.68 &  452.3 \\
CACTI            &  1000 &    180 & +22.94 &  458.8 \\
Lensless         &   300 &     45 &  +3.55 &  284.0 \\
MRI              &   150 &     30 & +48.25 & 5790.0 \\
SPC              &   200 &     15 & +12.21 & 2930.4 \\
CASSI (Alg~1)    &   500 &    300 &  +0.54 &    6.5 \\
CASSI (Alg~2)    &  2000 &   3200 &  +0.76 &    0.9 \\
Ptychography     &   800 &    420 &  +7.09 &   60.8 \\
\bottomrule
\end{tabular}
\end{table}

\subsection{RoIC Metric: Efficiency of Correction}

We define the \emph{Return on Invested Compute} (RoIC) as:
\begin{equation}
  \mathrm{RoIC}
  = \frac{\Delta\mathrm{PSNR}\;(\text{dB})}
         {T_{\mathrm{GPU}}\;(\text{hours})},
\end{equation}
where $T_{\mathrm{GPU}}$ is the wall-clock GPU time.

\noindent\textbf{Key observations:}
\begin{itemize}
  \item \textbf{MRI} achieves the highest RoIC (5790\,dB/GPU-hr) because coil
        sensitivity correction yields massive PSNR gains with a well-conditioned
        linear operator that converges rapidly.
  \item \textbf{Matrix/SPC} also show high RoIC due to the simplicity of the
        gain--bias correction model.
  \item \textbf{CASSI Algorithm~2} has the lowest RoIC (0.9\,dB/GPU-hr)
        because it jointly optimizes mask geometry and dispersion over 2000
        iterations, each requiring a full forward--adjoint cycle through the
        dispersive operator. Under the multi-parameter InverseNet mismatch,
        the absolute correction gains are modest ($+0.76$\,dB), reflecting
        the inherent difficulty of simultaneously correcting five coupled
        mismatch parameters.
  \item The practical implication is that Algorithm~1 (fixed dispersion,
        correct mask geometry only) is preferred when compute budget is
        limited, recovering 71\% of Algorithm~2's $\Delta$PSNR at
        $<$10\% of the cost.
\end{itemize}

\subsection{Scaling Considerations}

For large-scale deployment:
\begin{itemize}
  \item Correction is \emph{embarrassingly parallel} across scenes/slices ---
        multi-GPU scaling is near-linear.
  \item The OperatorGraph compilation overhead is amortized: once compiled,
        the same graph serves all scenes of a given modality.
  \item Memory footprint scales with the measurement dimension $m$, not the
        scene dimension $n$, making correction feasible even for
        high-resolution 3D volumes (e.g.\ $512^3$ CT).
\end{itemize}

\newpage

%% ============================================================
%% SUPPLEMENTARY NOTE 6: REAL-DATA VALIDATION
%% ============================================================
\section{Supplementary Note 6: Real-Data Validation Details}
\label{sec:realdata}

\subsection{CASSI Real Data: TSA Hyperspectral Camera}

The TSA (Thermal Snapshot Aperture) real dataset consists of 5 hyperspectral
scenes acquired with a coded aperture snapshot spectral imaging system.
Each scene has spatial resolution $660 \times 660$ with 28 spectral bands
(450--650\,nm) and a mask-shift step of 2 pixels per band.  The extended
measurement has width $660 + (28-1) \times 2 = 714$ pixels.

\paragraph{Experimental protocol.}
For each scene, we reconstruct with four methods: GAP-TV (classical
iterative), HDNet (mask-oblivious neural network), MST-S, and MST-L
(mask-guided spectral transformers).  Each method is applied under two
conditions: (i)~the calibrated mask provided with the dataset, and
(ii)~a perturbed mask with sub-pixel shift $dx = 0.5$\,px, $dy = 0.3$\,px.
Because the MST model has a hardcoded $256 \times 256$ spatial assumption,
the $660 \times 660$ real data is center-cropped to $256 \times 256$ for
MST reconstruction.  HDNet is fully convolutional and processes the full
spatial extent.

\paragraph{Metric.}
Because no ground-truth hyperspectral cube exists for the real scenes,
we use the measurement residual as a ground-truth-free diagnostic:
\begin{equation}
  r = \frac{\|\mathbf{y} - H\hat{\mathbf{x}}\|^2}{\|\mathbf{y}\|^2},
  \label{eq:meas_residual}
\end{equation}
where $\mathbf{y}$ is the measured snapshot and $\hat{\mathbf{x}}$ is the
reconstruction.  The residual ratio $r_{\mathrm{mismatch}} /
r_{\mathrm{calibrated}}$ quantifies the relative impact of operator
mismatch, independent of the (unknown) ground truth.

\begin{table}[htbp]
\centering
\caption{Supplementary Table S5: CASSI real-data measurement residuals.
Residual values $r$ (Eq.~\ref{eq:meas_residual}) and residual ratio
(mismatched/calibrated) across 5 TSA scenes and 4 methods.}
\label{tab:cassi_real}
\begin{tabular}{l l r r r}
\toprule
\textbf{Scene} & \textbf{Method} &
  \textbf{$r_{\mathrm{cal}}$} & \textbf{$r_{\mathrm{mis}}$} &
  \textbf{Ratio} \\
\midrule
Scene~1 & GAP-TV  & 0.00148 & 0.00298 & 2.0 \\
Scene~1 & HDNet   & 0.18506 & 0.18506 & 1.0 \\
Scene~1 & MST-S   & 0.09953 & 0.10746 & 1.1 \\
Scene~1 & MST-L   & 0.10055 & 0.09544 & 0.9 \\
\midrule
Scene~2 & GAP-TV  & 0.00210 & 0.00331 & 1.6 \\
Scene~2 & HDNet   & 0.14853 & 0.14853 & 1.0 \\
Scene~2 & MST-S   & 0.25357 & 0.23807 & 0.9 \\
Scene~2 & MST-L   & 0.25441 & 0.22811 & 0.9 \\
\midrule
Scene~3 & GAP-TV  & 0.00204 & 0.00320 & 1.6 \\
Scene~3 & HDNet   & 0.18291 & 0.18291 & 1.0 \\
Scene~3 & MST-S   & 0.09399 & 0.09160 & 1.0 \\
Scene~3 & MST-L   & 0.10127 & 0.09085 & 0.9 \\
\midrule
Scene~4 & GAP-TV  & 0.00148 & 0.00297 & 2.0 \\
Scene~4 & HDNet   & 0.14355 & 0.14355 & 1.0 \\
Scene~4 & MST-S   & 0.15586 & 0.14712 & 0.9 \\
Scene~4 & MST-L   & 0.17899 & 0.15893 & 0.9 \\
\midrule
Scene~5 & GAP-TV  & 0.00233 & 0.00421 & 1.8 \\
Scene~5 & HDNet   & 0.11325 & 0.11325 & 1.0 \\
Scene~5 & MST-S   & 0.11995 & 0.12004 & 1.0 \\
Scene~5 & MST-L   & 0.13008 & 0.12208 & 0.9 \\
\midrule
\textbf{Mean} & \textbf{GAP-TV} & \textbf{0.00189} & \textbf{0.00333} & \textbf{1.8} \\
\textbf{Mean} & \textbf{HDNet}  & \textbf{0.15466} & \textbf{0.15466} & \textbf{1.0} \\
\textbf{Mean} & \textbf{MST-S}  & \textbf{0.14458} & \textbf{0.14086} & \textbf{1.0} \\
\textbf{Mean} & \textbf{MST-L}  & \textbf{0.15306} & \textbf{0.13908} & \textbf{0.9} \\
\bottomrule
\end{tabular}
\end{table}

\noindent\textbf{Key findings.}
GAP-TV, which explicitly conditions on the coded aperture mask in its
forward--adjoint iterations, shows the expected sensitivity to mask mismatch
(mean ratio $1.8\times$).  HDNet, whose architecture processes the
measurement without explicit mask conditioning, is entirely insensitive
(ratio $1.0\times$).  The transformer methods (MST-S, MST-L) show ratios
near or below $1.0\times$ on real data---in sharp contrast to their severe
degradation in simulation ($-13.98$\,dB for MST-L).  This
simulation-to-hardware gap indicates that the real hardware mask already
contains uncorrected manufacturing imperfections, so the additional
$0.5$\,px perturbation is small relative to pre-existing errors.


\subsection{CACTI Real Data: Temporal Compressive Camera}

The CACTI real dataset consists of 4 dynamic scenes (duomino, hand,
pendulumBall, waterBalloon) acquired with a coded aperture compressive
temporal imaging system at $512 \times 512$ spatial resolution with
compression ratio 10 (10 video frames encoded per snapshot).

\begin{table}[htbp]
\centering
\caption{Supplementary Table S6: CACTI real-data measurement residuals
and total variation.  Residual ratio (mismatched/calibrated) and
total variation (TV) for 4 real scenes and 2 methods.}
\label{tab:cacti_real}
\begin{tabular}{l l r r r}
\toprule
\textbf{Scene} & \textbf{Method} &
  \textbf{$r_{\mathrm{cal}}$} & \textbf{$r_{\mathrm{mis}}$} &
  \textbf{Ratio} \\
\midrule
duomino       & GAP-TV     & $8.0\times10^{-6}$ & $8.5\times10^{-5}$ & 10.6 \\
duomino       & PnP-FFDNet & 0.00198             & 0.00403             &  2.0 \\
\midrule
hand          & GAP-TV     & $7.0\times10^{-6}$ & $7.7\times10^{-5}$ & 11.0 \\
hand          & PnP-FFDNet & 0.00249             & 0.00705             &  2.8 \\
\midrule
pendulumBall  & GAP-TV     & $3.7\times10^{-5}$ & $3.46\times10^{-4}$ &  9.4 \\
pendulumBall  & PnP-FFDNet & 0.00919             & 0.01150             &  1.3 \\
\midrule
waterBalloon  & GAP-TV     & $1.4\times10^{-5}$ & $1.47\times10^{-4}$ & 10.5 \\
waterBalloon  & PnP-FFDNet & 0.00264             & 0.00493             &  1.9 \\
\midrule
\textbf{Mean} & \textbf{GAP-TV}     & --- & --- & \textbf{10.4} \\
\textbf{Mean} & \textbf{PnP-FFDNet} & --- & --- & \textbf{2.0} \\
\bottomrule
\end{tabular}
\end{table}

\noindent\textbf{Key findings.}
GAP-TV shows an order-of-magnitude residual increase ($10.4\times$ mean)
under sub-pixel mask mismatch, consistent with the multiplicative error
amplification in temporal compression (a single mask error propagates
across all 10 compressed frames).  PnP-FFDNet shows moderate robustness
($2.0\times$), likely because the FFDNet denoiser provides an implicit
regularization that partially absorbs the mismatch artefacts.  The CACTI
results confirm that temporal compressive modalities are substantially
more sensitive to operator mismatch than spectral compressive modalities
(CASSI), a distinction that would be invisible without the unified
\textsc{Triad Decomposition} framework.


\subsection{Autonomous Calibration on Real Data}

\begin{table}[htbp]
\centering
\caption{Supplementary Table S7: Grid-search calibration results on
simulated data with known ground truth.  PSNR values in dB.}
\label{tab:scenario_iv}
\begin{tabular}{l r r r r r}
\toprule
\textbf{Modality} &
  \textbf{Sc.~II} & \textbf{Sc.~IV} & \textbf{Sc.~III} &
  \textbf{Time (s)} & \textbf{Recovery} \\
\midrule
CASSI  & 21.95 & 23.31 & 23.69 & 1234.3 & 78\% \\
CACTI  & 17.60 & 26.99 & 26.99 &   64.2 & 100\% \\
SPC    & 25.47 & 25.47 & 25.43 &  217.5 &  0\% \\
\bottomrule
\end{tabular}
\end{table}

\noindent CACTI achieves perfect calibration recovery because the mismatch
manifold is low-dimensional (2D spatial shift) and the temporal compression
provides high sensitivity to mask position.  CASSI achieves 78\% recovery;
the residual gap reflects the higher-dimensional mismatch space (mask
geometry plus dispersion) and pre-existing hardware imperfections.  SPC
gain-drift mismatch is not recoverable from the measurement residual alone,
because a multiplicative gain bias scales both signal and noise uniformly,
leaving the residual structure unchanged.  This demonstrates a fundamental
limitation of residual-based calibration for gain-type mismatches and
motivates the development of alternative diagnostics for \textsc{Gate~3}
in gain-dominated systems.


%% ============================================================
%% SUPPLEMENTARY TABLE S8: SSIM COMPARISON
%% ============================================================
\newpage
\section{Supplementary Table S8: SSIM Comparison Across Modalities}
\label{sec:ssim}

Table~\ref{tab:ssim} reports the structural similarity index (SSIM)
for all three validated photon-domain modalities under the three
primary scenarios. SSIM complements PSNR by capturing perceptual
quality, particularly structural distortions from operator mismatch.

\begin{table}[htbp]
\centering
\caption{SSIM comparison across modalities and scenarios.
Values are mean $\pm$ s.d.\ across test scenes/images.}
\label{tab:ssim}
\begin{tabular}{l l r r r}
\toprule
\textbf{Modality} & \textbf{Method} &
  \textbf{Sc.~I} & \textbf{Sc.~II} & \textbf{Sc.~III} \\
\midrule
\multicolumn{5}{l}{\emph{CASSI (10 KAIST scenes, 28 spectral bands)}} \\
\midrule
& GAP-TV   & $0.723 \pm 0.088$ & $0.612 \pm 0.084$ & $0.688 \pm 0.083$ \\
& HDNet    & $0.970 \pm 0.009$ & $0.756 \pm 0.074$ & $0.756 \pm 0.074$ \\
& MST-L    & $0.973 \pm 0.009$ & $0.744 \pm 0.069$ & $0.881 \pm 0.035$ \\
\midrule
\multicolumn{5}{l}{\emph{CACTI (6 benchmark videos, 8 temporal frames)}} \\
\midrule
& GAP-TV        & $0.848 \pm 0.083$ & $0.305 \pm 0.070$ & $0.794 \pm 0.063$ \\
& PnP-FFDNet    & $0.890 \pm 0.060$ & $0.216 \pm 0.076$ & $0.820 \pm 0.054$ \\
& ELP-Unfolding & $0.965 \pm 0.013$ & $0.308 \pm 0.076$ & $0.927 \pm 0.017$ \\
& EfficientSCI  & $0.973 \pm 0.012$ & $0.303 \pm 0.079$ & $0.927 \pm 0.017$ \\
\midrule
\multicolumn{5}{l}{\emph{SPC (11 standard images, 25\% compression)}} \\
\midrule
& FISTA-TV   & $0.852 \pm 0.046$ & $0.586 \pm 0.055$ & $0.759 \pm 0.036$ \\
& PnP-DRUNet & $0.880 \pm 0.047$ & $0.402 \pm 0.062$ & $0.645 \pm 0.050$ \\
& ISTA-Net   & $0.916 \pm 0.028$ & $0.584 \pm 0.077$ & $0.760 \pm 0.056$ \\
& HATNet     & $0.847 \pm 0.049$ & $0.648 \pm 0.091$ & $0.807 \pm 0.060$ \\
\bottomrule
\end{tabular}
\end{table}

\noindent\textbf{Key observations.}
The SSIM degradation pattern mirrors the PSNR findings: under Scenario~II
(mismatch), all methods collapse to a narrow quality range regardless of
their Scenario~I performance. In CASSI, the SSIM collapse is particularly
striking: MST-L drops from $0.973$ to $0.744$, a perceptual quality
loss that exceeds the improvement from a decade of solver development.
In CACTI, mismatch produces near-random SSIM values (${\sim}0.3$),
confirming that the reconstructed video frames bear little structural
resemblance to the ground truth under mismatch conditions.


%% ============================================================
%% SUPPLEMENTARY TABLE S9: CASSI SAM COMPARISON
%% ============================================================
\section{Supplementary Table S9: CASSI Spectral Angle Mapper (SAM)}
\label{sec:sam}

Table~\ref{tab:sam} reports the spectral angle mapper (SAM, in degrees)
for CASSI, which measures spectral fidelity independently of intensity.
Lower values indicate better spectral reconstruction.

\begin{table}[htbp]
\centering
\caption{CASSI spectral angle mapper (SAM, degrees). Lower is better.
Values are mean $\pm$ s.d.\ across 10 KAIST scenes.}
\label{tab:sam}
\begin{tabular}{l r r r}
\toprule
\textbf{Method} &
  \textbf{Sc.~I} & \textbf{Sc.~II} & \textbf{Sc.~III} \\
\midrule
GAP-TV      & $16.66 \pm 3.32$ & $24.27 \pm 2.92$ & $25.97 \pm 2.54$ \\
PnP-HSI-CNN & $16.10 \pm 3.25$ & $23.73 \pm 2.81$ & $18.66 \pm 2.75$ \\
HDNet       & $6.67 \pm 0.96$  & $17.03 \pm 2.76$ & $17.03 \pm 2.76$ \\
MST-L       & $7.44 \pm 1.24$  & $23.92 \pm 4.58$ & $11.74 \pm 1.32$ \\
\bottomrule
\end{tabular}
\end{table}

\noindent\textbf{Key observations.}
Spectral fidelity degrades severely under mismatch: MST-L's SAM increases
from $7.44^\circ$ (near-perfect spectral reconstruction) to $23.92^\circ$
(substantial spectral mixing). Oracle correction partially recovers
spectral quality ($11.74^\circ$), but the residual SAM error confirms
that multi-parameter mismatch produces irreversible spectral distortions
not fully correctable by mask-only correction. The SAM metric is
particularly valuable for remote sensing applications where spectral
accuracy is more important than spatial quality.


%% ============================================================
%% SUPPLEMENTARY TABLES S10--S11: GATE 1 & GATE 2 VALIDATION
%% ============================================================
\newpage
\section{Supplementary Tables S10--S11: Gate~1 and Gate~2 Validation}
\label{sec:gate12}

The main text argues that Gate~3 (operator mismatch) is the dominant
bottleneck under standard operating conditions. To demonstrate that
Gate~1 (information deficiency) and Gate~2 (carrier budget) are real and
measurable, we sweep the compression level and noise level across all
seven validated modalities while keeping the forward model perfectly
calibrated.  All experiments use classical solvers (no deep-learned
components) to ensure the results reflect information-theoretic limits
rather than solver-specific behaviour.

\subsection{Gate~1: Information Deficiency (Extreme Compression)}

Table~\ref{tab:gate1} reports reconstruction PSNR (dB) as the
compression ratio / sampling rate / blur width is driven to extreme
values.  The forward model is ideal in every case; any PSNR loss is
attributable solely to information lost in the measurement process.

\begin{table}[htbp]
\centering
\caption{Gate~1 validation: PSNR (dB) under extreme compression for 7
modalities.  ``Best solver'' reports the better of the two classical
methods tested.  Bold entries highlight the regime where PSNR falls
below $15$\,dB (severe information loss).}
\label{tab:gate1}
\small
\begin{tabular}{l l r r r r r}
\toprule
\textbf{Modality} & \textbf{Parameter} &
  \multicolumn{5}{c}{\textbf{PSNR (dB) at sweep points}} \\
\cmidrule(lr){3-7}
 & & \textbf{Nominal} & \textbf{Moderate} & \textbf{Severe} & \textbf{Extreme} & \textbf{Critical} \\
\midrule
CACTI & CR
  & 25.9 \scriptsize{(8)} & 24.1 \scriptsize{(16)} & 22.7 \scriptsize{(32)} & 20.6 \scriptsize{(64)} & --- \\
CASSI & Transmittance
  & 26.1 \scriptsize{(50\%)} & 27.7 \scriptsize{(25\%)} & 27.7 \scriptsize{(10\%)} & 27.6 \scriptsize{(5\%)} & 27.4 \scriptsize{(2\%)} \\
SPC & CS ratio
  & 28.3 \scriptsize{(25\%)} & 24.0 \scriptsize{(10\%)} & 21.2 \scriptsize{(5\%)} & 16.2 \scriptsize{(2\%)} & \textbf{14.4} \scriptsize{(1\%)} \\
MRI & Sampling rate
  & 28.8 \scriptsize{(25\%)} & 25.1 \scriptsize{(10\%)} & 24.3 \scriptsize{(5\%)} & 23.9 \scriptsize{(2\%)} & --- \\
CT & Projection angles
  & 22.1 \scriptsize{(180)} & 22.0 \scriptsize{(90)} & 21.2 \scriptsize{(30)} & 20.7 \scriptsize{(10)} & 17.8 \scriptsize{(5)} \\
Lensless & Blur $\sigma$ (px)
  & 36.6 \scriptsize{(1)} & 25.5 \scriptsize{(3)} & 23.1 \scriptsize{(5)} & 20.5 \scriptsize{(10)} & 17.9 \scriptsize{(20)} \\
Ptychography & Scan positions
  & 15.0 \scriptsize{(16)} & 12.4 \scriptsize{(9)} & 12.0 \scriptsize{(4)} & 9.9 \scriptsize{(1)} & --- \\
\bottomrule
\end{tabular}
\end{table}

\noindent\textbf{Key observations.}
\begin{itemize}
  \item \textbf{SPC} and \textbf{Lensless} show the most dramatic Gate~1 collapse.
    SPC PSNR drops $13.9$\,dB from $25$\% to $1$\% CS ratio; lensless drops
    $18.7$\,dB as the PSF blur widens from $\sigma{=}1$ to $\sigma{=}20$\,px.
    In both cases the loss is irrecoverable: no solver can reconstruct
    information that was never measured.
  \item \textbf{CASSI} is a notable exception: reducing mask transmittance from
    $50$\% to $2$\% produces no degradation (mean PSNR actually increases by
    ${\sim}1.5$\,dB). This occurs because sparser coded apertures reduce
    spectral mixing in the multiplexed measurement, making the
    demixing problem easier for iterative solvers.  The CASSI null space
    is effectively unchanged because each spectral band still contributes
    signal through the remaining open mask pixels.
  \item \textbf{CACTI} degrades monotonically with compression ratio
    ($25.9 \to 20.6$\,dB over $8\times$ CR increase), consistent with the
    growing null space when fewer temporal measurements are available.
  \item \textbf{CT} shows SART outperforming FBP at sparse angles
    ($21.2$\,dB vs.\ $15.8$\,dB at 10 angles), illustrating that iterative
    solvers can partially compensate for angular undersampling via their
    implicit prior, but ultimately both methods converge at $\leq 5$ angles.
\end{itemize}


\subsection{Gate~2: Carrier Budget (Noise Sweep)}

Table~\ref{tab:gate2} reports reconstruction PSNR as the photon budget
or noise level is swept to extreme values while the compression is held
at a well-posed nominal value and the forward model is ideal.

\begin{table}[htbp]
\centering
\caption{Gate~2 validation: PSNR (dB) under increasing noise for 7
modalities.  ``Best solver'' reports the better of the two classical
methods tested.  Noise parameter: photon count for Poisson-dominated
modalities, Gaussian $\sigma$ (fraction of signal range) for
MRI and SPC.  Bold entries indicate PSNR ${<}\,15$\,dB.}
\label{tab:gate2}
\small
\begin{tabular}{l l r r r r r}
\toprule
\textbf{Modality} & \textbf{Noise param.} &
  \multicolumn{5}{c}{\textbf{PSNR (dB) at sweep points}} \\
\cmidrule(lr){3-7}
 & & \textbf{Low noise} & & \textbf{Moderate} & & \textbf{Extreme} \\
\midrule
CACTI & Photon count
  & 24.8 \scriptsize{(10k)} & 23.3 \scriptsize{(1k)} & 18.4 \scriptsize{(100)} & \textbf{10.5} \scriptsize{(10)} & --- \\
CASSI & Photon count
  & 24.6 \scriptsize{(10k)} & 22.3 \scriptsize{(1k)} & 20.2 \scriptsize{(100)} & 17.0 \scriptsize{(10)} & --- \\
SPC & $\sigma$ (frac.)
  & 27.9 \scriptsize{(0)} & 27.8 \scriptsize{(.01)} & 25.3 \scriptsize{(.05)} & 21.6 \scriptsize{(.1)} & \textbf{13.5} \scriptsize{(.3)} \\
MRI & $\sigma$ (frac.)
  & 28.8 \scriptsize{(0)} & 17.1 \scriptsize{(.01)} & \textbf{12.8} \scriptsize{(.05)} & \textbf{12.1} \scriptsize{(.1)} & \textbf{11.0} \scriptsize{(.3)} \\
CT & Photon count
  & 22.1 \scriptsize{(100k)} & 21.9 \scriptsize{(10k)} & 21.0 \scriptsize{(1k)} & 18.3 \scriptsize{(100)} & --- \\
Lensless & Photon count
  & 40.9 \scriptsize{(10k)} & 32.8 \scriptsize{(1k)} & 23.1 \scriptsize{(100)} & \textbf{13.6} \scriptsize{(10)} & --- \\
Ptychography & Photon count
  & 13.5 \scriptsize{(10k)} & \textbf{10.3} \scriptsize{(1k)} & \textbf{10.0} \scriptsize{(100)} & \textbf{9.7} \scriptsize{(10)} & --- \\
\bottomrule
\end{tabular}
\end{table}

\noindent\textbf{Key observations.}
\begin{itemize}
  \item Every modality exhibits monotonic PSNR degradation with
    increasing noise, confirming that Gate~2 failures are universal.
  \item \textbf{MRI} is the most noise-sensitive modality: adding just
    $\sigma{=}0.01$ Gaussian noise to k-space drops CS-wavelet PSNR from
    $28.8$\,dB to $17.1$\,dB, a $11.7$\,dB cliff. This is because MRI's
    Fourier encoding concentrates signal energy in a few k-space samples,
    so additive noise corrupts the information-dense center region
    disproportionately.
  \item \textbf{Lensless} imaging shows the widest Gate~2 dynamic range:
    $40.9$\,dB at 10,000 photons to $13.6$\,dB at 10 photons, a $27.3$\,dB
    span. This reflects the well-conditioned forward operator (diffuser PSF),
    which faithfully propagates both signal and noise.
  \item \textbf{CT} is the most noise-robust modality ($3.8$\,dB drop from
    $10^5$ to $10^2$ photons), because the many-angle sinogram provides
    substantial redundancy that averages out Poisson noise.
  \item The cliff-edge nature of Gate~2 is most dramatic in CACTI
    ($14.3$\,dB drop from $10{,}000$ to $10$ photons) and SPC
    ($14.4$\,dB drop from $\sigma{=}0$ to $\sigma{=}0.3$), where the
    compressed measurement offers no noise-averaging redundancy.
\end{itemize}

\noindent\textbf{Comparison with Gate~3.}
Under standard operating conditions (adequate compression, adequate
photon budget), the Gate~3 (mismatch) degradation reported in
Supplementary Table~S1 is typically $3$--$20$\,dB---comparable in
magnitude to Gate~1 and Gate~2 degradation, but \emph{correctable} by
operator refinement.  This is the key distinction: Gate~1 and Gate~2
losses are information-theoretic and irrecoverable, whereas Gate~3
losses are recoverable via calibration.  The practical dominance of
Gate~3 arises because modern instruments operate well above the Gate~1
and Gate~2 thresholds, leaving operator mismatch as the binding
constraint.


%% ============================================================
%% SUPPLEMENTARY NOTE 7: CLINICAL CT QC VALIDATION
%% ============================================================
\newpage
\section{Supplementary Note 7: Clinical CT Quality Assurance Validation}
\label{sec:ctqc}

The \textsc{Triad Decomposition} framework has been translated to clinical
CT quality assurance (QA) through the CT QC Copilot module, which maps
the three gates to clinical failure modes and implements ACR CT
accreditation standards.

\subsection{Gate Mapping to Clinical Failure Modes}

\begin{table}[htbp]
\centering
\caption{Mapping of Triad Decomposition gates to clinical CT QC failure modes.}
\label{tab:gate_mapping}
\begin{tabular}{l l l}
\toprule
\textbf{Gate} & \textbf{Research Interpretation} & \textbf{Clinical Interpretation} \\
\midrule
Gate~1 & Information deficiency       & Protocol design inadequacy \\
       & (null space, compression)     & (insufficient projections, FOV) \\
\midrule
Gate~2 & Carrier budget               & Dose budget \\
       & (SNR, photon count)          & (noise floor vs.\ diagnostic need) \\
\midrule
Gate~3 & Operator mismatch            & Scanner calibration drift \\
       & (mask shift, PSF error)      & (HU drift, CoR offset, gain) \\
\bottomrule
\end{tabular}
\end{table}

\subsection{ACR Metric Validation}

Nine ACR-aligned QC metrics are computed automatically from CT phantom
(Gammex 464) scans.  Table~\ref{tab:acr_metrics} reports the measurement
agreement between the CT QC Copilot and console-reported values on a GE
Revolution Apex scanner (120~kVp, 200~mAs, 5~mm axial).

\begin{table}[htbp]
\centering
\caption{ACR metric agreement: CT QC Copilot vs.\ console gold standard.}
\label{tab:acr_metrics}
\begin{tabular}{l r r r r}
\toprule
\textbf{Metric} & \textbf{Copilot} & \textbf{Console} &
  \textbf{Deviation} & \textbf{\% Tolerance} \\
\midrule
CT\# Water (HU)          & 0.3     & 0.2     & 0.1   & 2\% \\
CT\# Bone (HU)           & 952     & 953     & 1.0   & 5\% \\
CT\# Air (HU)            & $-$997  & $-$998  & 1.0   & 5\% \\
CT\# Acrylic (HU)        & 120     & 121     & 1.0   & 7\% \\
CT\# Polyethylene (HU)   & $-$96   & $-$96   & 0.0   & 0\% \\
Geometric accuracy (mm)  & 0.15    & 0.10    & 0.05  & 2.5\% \\
Slice thickness (mm)     & 5.08    & 5.02    & 0.06  & 4\% \\
Uniformity (HU)          & 1.2     & 1.1     & 0.1   & 2\% \\
Noise std.\ dev.\ (HU)  & 4.6     & 4.5     & 0.1   & 10\% \\
Spatial resolution (lp/cm) & 6.2   & 6.0     & 0.2   & 4\% \\
\bottomrule
\end{tabular}
\end{table}

\noindent All metrics are within $\leq 10$\% of clinical tolerance bands,
confirming that automated computation agrees with manual physicist
measurements to within acceptable precision.

\subsection{Drift Detection Performance}

Statistical process control (SPC) using five Western Electric rules was
evaluated on a 30-scanner simulated fleet over 12 monthly measurement
cycles.  Four scanners were programmed with known drift trajectories
(noise drift: 2 scanners; uniformity drift: 1 scanner; CT number drift:
1 scanner).

\begin{itemize}
  \item \textbf{Sensitivity}: 100\% (4/4 drifting scanners detected;
    95\% CI: 39.8\%--100\%)
  \item \textbf{Specificity}: 100\% (0/26 false positives on stable
    scanners; 95\% CI: 86.8\%--100\%)
  \item \textbf{Early warning lead time}: 3--6 months before ACR
    threshold exceedance
\end{itemize}

\noindent The automated detection provides a proactive maintenance window
that is not achievable with manual threshold-based QA, which only
triggers at the point of failure.

\subsection{Workflow Efficiency}

\begin{table}[htbp]
\centering
\caption{Per-scanner QC workflow time comparison.}
\label{tab:workflow}
\begin{tabular}{l r r}
\toprule
\textbf{Stage} & \textbf{Manual (min)} & \textbf{Copilot (min)} \\
\midrule
DICOM ingestion          & 5.0   & 0.01 \\
Metric computation       & 25.0  & 0.01 \\
Threshold evaluation     & 5.0   & $<$0.01 \\
Trend analysis           & 10.0  & $<$0.01 \\
Report generation        & 15.0  & 0.02 \\
Physicist review/sign-off & 7.0  & 4.0 \\
\midrule
\textbf{Total}           & $67 \pm 12$ & $4.2 \pm 0.8$ \\
\bottomrule
\end{tabular}
\end{table}

\noindent The 94\% reduction in per-scanner QC time (from $67 \pm 12$ to
$4.2 \pm 0.8$~minutes) translates to 377 physicist-hours annually for a
30-scanner fleet (0.18~FTE reallocation).  Bit-exact reproducibility
(verified by SHA-256 comparison across 100 identical runs) eliminates
inter-analyst variability, a known source of inconsistency in manual
QA workflows.

\bibliographystyle{unsrtnat}
\bibliography{pwm_flagship}

\end{document}
